% ============================================================================
% SECTION 1: INTRODUCTION
% ============================================================================
\section{Introduction}\label{sec:intro}

\subsection{Motivation}

Side-channel attacks exploit implementation leakage rather than weaknesses in
the abstract cipher alone.  In many attacks on lightweight block ciphers, the
local statistical object is a keyed S-box computation: an adversary observes or
controls an input nibble or byte, measures a noisy leakage trace, and compares
key hypotheses through predicted intermediate values.  Correlation power
analysis (CPA) is the standard example: each trace updates the adversary's
belief about the key using a leakage model such as Hamming weight.

The exact distribution induced by a keyed bijective S-box is not a smooth
interior point of a probability simplex.  For an $s$-bit bijective S-box
$S:\{0,1\}^s\to\{0,1\}^s$ and key hypothesis $k\in\{0,1\}^s$, the pair
$(a,b)$ with uniform $a$ and $b=S(a\oplus k)$ is supported on only $2^s$ points
of the $2^{2s}$-point input-output space.  Hence $P_k$ is a singular boundary
distribution.  This matters because ordinary Fisher--Rao tensors, geodesics,
and dual-flat coordinates are defined on full-support statistical models, not
directly at support-deficient endpoints.

This paper develops a boundary information-geometric framework for this
side-channel setting.  The central discipline is to keep separate three
objects that are often conflated in informal geometric descriptions:
\begin{enumerate}[label=\normalfont(\roman*)]
  \item the regular full-support Walsh exponential family, where ordinary
  Fisher geometry is valid;
  \item the finite boundary orbit of exact key-induced S-box distributions;
  \item the adversary's posterior belief simplex over key hypotheses, where
  CPA updates take place.
\end{enumerate}
This separation makes it possible to use information geometry in a controlled
way: boundary distributions are represented through covariance limits, whereas
Bayesian key recovery remains a regular trajectory on the belief simplex as
long as the prior has full support and the leakage likelihood is positive.

\begin{figure}[t]
\centering
\resizebox{\textwidth}{!}{%
\begin{tikzpicture}[
  font=\small,
  node distance=11mm and 12mm,
  box/.style={
    draw,
    rounded corners=1.5pt,
    align=center,
    inner sep=5pt,
    minimum width=34mm,
    minimum height=11mm
  },
  wide/.style={
    draw,
    rounded corners=1.5pt,
    align=center,
    inner sep=5pt,
    minimum width=46mm,
    minimum height=12mm
  },
  arr/.style={-{Latex[length=2mm]},thick}
]
\node[wide] (sbox) {Keyed S-box law\\$a\mapsto b=S(a\oplus k)$};
\node[box, below left=of sbox] (boundary) {Boundary orbit\\$\{P_k:k\in\{0,1\}^s\}$};
\node[box, below=of sbox] (interior) {Regular Walsh\\exponential family};
\node[box, below right=of sbox] (belief) {Key-belief simplex\\$q_i(k)$};
\node[wide, below=of interior] (diagnostics) {Security diagnostics\\active Fisher structure, CPA gaps};

\draw[arr] (sbox) -- node[left,align=center] {exact graph\\distribution} (boundary);
\draw[arr] (sbox) -- node[right,align=center] {leakage\\likelihoods} (belief);
\draw[arr] (boundary) -- node[above,align=center] {Walsh covariance\\limits} (interior);
\draw[arr] (interior) -- node[left,align=center] {active\\eigenspace} (diagnostics);
\draw[arr] (belief) -- node[right,align=center] {Bayesian\\updates} (diagnostics);
\end{tikzpicture}
}
\caption{Boundary-to-belief framework. The exact keyed S-box laws lie on a
singular boundary orbit, the Fisher calculations are made in a regular Walsh
exponential family or in boundary covariance limits, and CPA updates evolve on
the separate key-belief simplex.}
\label{fig:framework}
\end{figure}

\subsection{Relation to prior work}

The information-geometric background is the Fisher--Rao and dual-affine
geometry of finite statistical models and regular exponential families
\cite{amari2000methods,ay2017information,chentsov1982statistical}.  The
closest coding-theoretic precedent is the work of Ikeda, Tanaka, and Amari on
LDPC and turbo codes~\cite{ikeda2004information}: parity-check statistics
define a regular information-geometric model and decoding trajectories are
interpreted through projections.  Here, the sufficient statistics are
Walsh--Hadamard characters on S-box input-output pairs, and the inference
trajectory is the CPA posterior over key hypotheses.

On the cryptographic side, the work is related to classical Walsh/LAT and DDT
analysis of vectorial Boolean functions and S-boxes, including Nyberg's
differential uniformity~\cite{nyberg1994differentially}.  Four-bit S-boxes and
optimal 4-bit permutations have been classified in the Boolean-function
literature, notably by Leander and Poschmann~\cite{leander2007classification}.
The present paper does not propose a new classification of 4-bit S-boxes and
does not claim that standard S-boxes with the same Walsh spectrum can be
distinguished by spectral invariants.  Instead, it asks a different question:
what is the correct information-geometric object associated with exact keyed
S-box distributions, and how should that object be connected to side-channel
belief updates?

On the side-channel side, the leakage model and posterior update are consistent
with differential and correlation power analysis
\cite{kocher1999differential,brier2004correlation} and with
information-theoretic analyses of key recovery attacks
\cite{standaert2009unified,veyrat2009adaptive}.  The present framework is not a
new attack and not a countermeasure.  It is a formal model for separating
baseline invariants forced by bijectivity from pairwise leakage gaps that
control key-hypothesis discrimination.

Table~\ref{tab:sota-positioning} summarizes the intended position of the
framework relative to standard S-box and side-channel criteria.

\begin{table}[t]
\centering
\caption{Relation to existing metrics.}
\label{tab:sota-positioning}
\footnotesize
\setlength{\tabcolsep}{3pt}
\begin{tabular}{@{}p{0.22\textwidth}p{0.34\textwidth}p{0.34\textwidth}@{}}
\toprule
Metric or framework & Standard role & Role in this paper \\
\midrule
Walsh/LAT spectrum & Measures linear correlations of vectorial Boolean
functions and supports nonlinearity analysis. & Provides coordinates for the
key-induced boundary orbit and for closed-form covariance diagnostics. \\
DDT and differential uniformity & Measures differential propagation and
classical S-box resistance to differential attacks~\cite{nyberg1994differentially}. &
Used as part of the surrounding S-box context, but not replaced or re-ranked by
the proposed geometry. \\
4-bit S-box classification & Classifies small permutations up to established
equivalences~\cite{leander2007classification}. & Treated as prior knowledge;
the paper does not claim a new classification or a new optimality criterion. \\
CPA and DPA likelihood models & Explain key recovery from noisy implementation
leakage~\cite{kocher1999differential,brier2004correlation}. & Interpreted as
Bayesian dynamics on the key-belief simplex with pairwise likelihood gaps. \\
Mutual information and key-rank measures & Quantify leakage and attack success
at the experiment or trace-distribution level~\cite{standaert2009unified,veyrat2009adaptive}. &
Complemented by local structural diagnostics for keyed S-box laws and
posterior-odds growth. \\
Information geometry of codes & Studies decoding and projections in regular
statistical models~\cite{ikeda2004information}. & Adapted to singular S-box
graph distributions by separating boundary covariance, regular interiors, and
belief-simplex updates. \\
\bottomrule
\end{tabular}
\end{table}

\subsection{Setting and notation}

Throughout, $s\geq1$ denotes the S-box width in bits,
$S:\B{s}\to\B{s}$ is a fixed bijection, $k\in\B{s}$ is a key hypothesis, and
$a\in\B{s}$ is a uniformly distributed plaintext nibble or byte.  All additions
inside bit strings are over $\F_2$.  The inner product is
$\alpha\cdot a=\bigoplus_{i=1}^s \alpha_i a_i$.  For
$(\alpha,\beta)\in\B{s}\times\B{s}$, the Walsh character is
\[
  \chi_{\alpha\beta}(a,b)=(-1)^{\alpha\cdot a\oplus\beta\cdot b}.
\]
We write $N=2^s$ and $\mathcal X=\B{s}\times\B{s}$.

\subsection{Contributions}\label{subsec:contributions}

\begin{enumerate}[label=\textbf{C\arabic*.}]

\item \textbf{Boundary model for key-induced S-box distributions.}
We model the exact keyed S-box law
\[
  P_k(a,b)=2^{-s}\mathbf 1[b=S(a\oplus k)]
\]
as a boundary distribution on the input-output simplex.  The Walsh translation
rule
\[
  \widehat P_k(\alpha,\beta)
  =(-1)^{\alpha\cdot k}\widehat P_0(\alpha,\beta)
\]
is used as a coordinate identity for the finite key orbit, not as a standalone
cryptanalytic novelty.

\item \textbf{Closed-form boundary covariance and active Fisher structure.}
We derive the boundary covariance form
\[
  \mathcal I^{\partial}_{ij}(k)=\operatorname{Cov}_{P_k}(T_i,T_j)
\]
in Walsh coordinates and prove that, for every bijection $S:\B{s}\to\B{s}$,
this matrix has rank $2^s-1$ and active eigenvalues all equal to $2^s$ in the
unnormalised Walsh basis.  Thus the active boundary geometry is a structural
property of bijective S-box graphs, rather than an artefact of the three
illustrative 4-bit examples.

\item \textbf{Regular interior and belief-simplex geometry.}
We explicitly separate the regular Walsh exponential family, where the usual
Amari dual-flat structure and Pythagorean identities apply, from the singular
boundary orbit.  The adversary's key belief is treated as a separate full-
support simplex, equipped with its own Fisher--Rao metric.

\item \textbf{CPA as Bayesian belief dynamics with pairwise leakage gaps.}
Under Gaussian leakage, each CPA trace is represented as a Bayesian update on
the key-belief simplex.  The Hamming-weight variance
$s/(4\sigma^2)$ is identified as a baseline forced by bijectivity and uniform
inputs.  The quantities that distinguish key hypotheses are the pairwise
separations
\[
  \Delta(k,k^*)=\frac1{2^s}\sum_a
  \bigl(f(S(a\oplus k^*))-f(S(a\oplus k))\bigr)^2,
\]
which determine expected log-Bayes gains.

\item \textbf{Reproducible diagnostics on lightweight S-boxes.}
For PRESENT, GIFT, and SKINNY, we compute Walsh spectra, boundary covariance
matrices, local spectral separation proxies, pairwise Hamming-weight leakage
gaps, and the projected active-eigenspace curvature diagnostic.  The examples
are used to separate invariants forced by bijectivity or by the shared Walsh
amplitude distribution from quantities that can vary across S-boxes.
\end{enumerate}

\subsection{Scope}\label{subsec:scope}

The paper focuses on isolated bijective S-boxes under uniform inputs and on an
additive Gaussian leakage model for CPA belief updates.  It does not address
complete multi-round cipher security, masking, higher-order leakage,
non-Gaussian leakage, profiling attacks, or implementation-specific trace
effects.  Those topics require additional models beyond the boundary
distribution and belief-update framework developed here.
